% Cote 119: offprints, reports, courses, CV and letters (typed and printed part): archive cote 119, site page 137 (PDF page 138).
% typescript. Draft compared: none.
% AI-transcribed by Claude Opus 5.5 (claude-opus-5-5, adaptive thinking, effort high), Message Batches API, batch r119_full (msgbatch_01DwZNm5Y7YsKBWyAN4sFs8w), 2026-11-13; not yet checked by eye.
\scanpage{119/137}
\uline{Rapport d'activité} (1.10.1984 - 15.12.1984)

par Alexandre Grothendieck, attaché de recherches au CNRS

Les mois écoulés ont été consacrés avant tout à la préparation des
volumes 1 et 2 des Réflexions Mathématiques, qui paraîtront sans doute
courant 1985, sous les titres ``\uline{Récoltes et Semailles}'' et ``\uline{Pursuing}
\uline{Stacks,} part 1 : \uline{The Modelizing Story}''. Je prévois que cette préparati-%
on m'absorbera jusque vers la fin février, après quoi je compte repren-%
dre la réflexion sur les fondements d'une ``algèbre topologique''\hand{(}\unsure{commencée|commencé}
avec la première partie de ``A la Poursuite des Champs''\hand{)}\unsure{,} en reprenant
le fil de la réflexion de l' ``Histoire de Modèles'' là où je l'avais
laissé. Je donne quelques précisions au sujet de ce programme de tra-%
vail dans ``Esquisse d'un Programme'', \uline{par.7.}\hand{]} Il s'agit ici de refondre,
dans une discipline autonome (que je propose d'appeler ``\uline{algèbre topologi-}%
\uline{que}'') qui jouerait le rôle un peu du pendant ``algébrique'' de la topologie
générale, un certain nombre de visions parcellaires provenant de l'al-%
gèbre (co)homologique (commutative ou non-commutative), l'algèbre homo-%
topique, le formalisme algébrico-géométrique des topos (qui est une sor-%
te de paraphrase dûment généralisée de la topologie générale), y compris
le formalisme des champs, gerbes etc développé dans la thèse de Giraud,
\inserted{enfin} la théorie des n-catégories\hand{,} et \inserted{celle} des n-champs de telles n-catégories
(encore dans les limbes). Le besoin d'une telle discipline nouvelle,
et certaines des principales idées-force que je compte développer dans
\xout{ilxsqui} le maître d'oeuvre ``A la P\inserted{o}ursuite des Champs'',me sont apparus
progressivement entre 1955 et la fin des années soixante, en contre-%
point \xout{constant} ininterrompu avec le développement d'une ``\uline{géométrie}
\uline{arithmétique}'', synthèse (entièrement imprévue encore jusqu'aux débuts
des années soixante) de la \unsure{géométrié|géométrie} algébrique, \inserted{de} la topologie et \inserted{de} l'arith-%
métique. Je développe \xout{\illegible{1}}des réflexions rétrospectives au sujet de
la naissance, l'essor et la fortune (et parfois, infortune) ultérieure
de ces disciplines nouvelles, ici et là au cours de la réflexion plus
vaste poursuivie dans ``Récoltes et Semailles''. Qu'il me suffise ici de
dire que dans mes propres motivations, aujourd'hui comme il y a vingt
ans, l'algèbre topologique (laissée pour compte \deleted{depuis} \inserted{après} mon ``départ''
en 1970) est avant tout un des principaux \uline{outils d'appoint} pour le
développement de cette ``géométrie arithmétique'', dont le développement
jusqu'au stade d'une pleine maturité m'apparaît comme une des \uline{principa-}%
les tâches qui se posent à la mathématique de notre \uline{temps.}\hand{]} Un des sig-%
nes principaux d'une telle maturité serait une maîtrise complète des

\mspage{137}
